\begin{lemma}
\label{lemma-upper-shriek-finite}
Let $\varphi : R \to A$ be a finite map of Noetherian rings.
Then $\varphi^!$ is isomorphic to the functor
$R\Hom(A, -) : D(R) \to D(A)$ from
Section \ref{section-trivial}.
\end{lemma}

\begin{proof}
Suppose that $A$ is generated by $n > 1$ elements over $R$.
Then can factor $R \to A$ as a composition of two finite ring maps
where in both steps the number of generators is $< n$.
Since we have Lemma \ref{lemma-composition-shriek-algebraic} and
Lemma \ref{lemma-composition-right-adjoints}
we conclude that it suffices
to prove the lemma when $A$ is generated by one element over $R$.
Since $A$ is finite over $R$, it follows that $A$ is a quotient
of $B = R[x]/(f)$ where $f$ is a monic polynomial in $x$
(Algebra, Lemma \ref{algebra-lemma-finite-is-integral}).
Again using the lemmas on composition and the fact that we
have agreement for surjections by definition, we conclude that
it suffices to prove the lemma for $R \to B = R[x]/(f)$.
In this case, the functor $\varphi^!$ is isomorphic to
$K \mapsto K \otimes_R^\mathbf{L} B$; you prove this by
using Lemma \ref{lemma-compute-for-effective-Cartier-algebraic}
for the map $R[x] \to B$ (note that the shift in the definition
of $\varphi^!$ and in the lemma add up to zero).
For the functor $R\Hom(B, -) : D(R) \to D(B)$ we can use
Lemma \ref{lemma-RHom-is-tensor-special}
to see that it suffices to show $\Hom_R(B, R) \cong B$
as $B$-modules. Suppose that $f$ has degree $d$.
Then an $R$-basis for $B$ is given by $1, x, \ldots, x^{d - 1}$.
Let $\delta_i : B \to R$, $i = 0, \ldots, d - 1$
be the $R$-linear map which picks off the coefficient
of $x^i$ with respect to the given basis. Then
$\delta_0, \ldots, \delta_{d - 1}$ is a basis for $\Hom_R(B, R)$.
Finally, for $0 \leq i \leq d - 1$ a computation shows that
$$
x^i \delta_{d - 1} =
\delta_{d - 1 - i} + b_1 \delta_{d - i} + \ldots + b_i \delta_{d - 1}
$$
for some $c_1, \ldots, c_d \in R$\footnote{If
$f = x^d + a_1 x^{d - 1} + \ldots + a_d$, then
$c_1 = -a_1$, $c_2 = a_1^2 - a_2$, $c_3 = -a_1^3 + 2a_1a_2 -a_3$, etc.}.
Hence $\Hom_R(B, R)$ is a principal $B$-module with generator
$\delta_{d - 1}$. By looking
at ranks we conclude that it is a rank $1$ free $B$-module.
\end{proof}